\documentclass[12pt]{article}

\PassOptionsToPackage{table}{xcolor}
\PassOptionsToPackage{a5paper, landscape, margin=1cm}{geometry}

\usepackage{
	fontspec,
	geometry,
	parskip,
	float,
	wrapfig,
	xcolor,
	tikz,
	amsmath,
	newtxmath
}

\def\hl{\colorbox{yellow}}
\def\hlr{\rowcolor{yellow}}
\def\True{\mathrm{True}}
\def\False{\mathrm{False}}

\setmainfont{Times New Roman}
\pagestyle{empty}

\title{Calculating Metric Dimension with Boolean Satisfiability}
\author{https://github.com/wijayarakhmat98/metric-dimension}

\begin{document}

\maketitle
\thispagestyle{empty}

\newpage
\section*{Boolean Satisfiability}

Boolean satisfiability is the process of finding boolean values for the variables in a boolean expression that makes the expression true.

For example, consider the following expression.

\[ (x_1 \land x_2) \lor (x_1 \land x_3) \]

A truth table can be constructed for the expression.

\[
\begin{array}{cccc}
	   x_1 &    x_2 &    x_3 & (x_1 \land x_2) \lor (x_1 \land x_3) \\
	\\
	\False & \False & \False & \False \\
	 \True & \False & \False & \False \\
	\False &  \True & \False & \False \\
	 \True &  \True & \False &  \True \\
	\False & \False &  \True & \False \\
	 \True & \False &  \True & \False \\
	\False &  \True &  \True &  \True \\
	 \True &  \True &  \True &  \True
\end{array}
\]

Such, $x_1 = \False$, $x_2 = \True$, $x_3 = \True$ is one solution that satisfies the expression.

\newpage
\section*{Boolean Satisfiability}

Boolean satisfiability does not specify which solution it will return. Rather, it simply finds one, called a model.

When a solution exists, the expression is said to be satisfiable. When no solution exists, the expression is said to be unsatisfiable.

For the notation in boolean satifiability, rather than using $\land\;\lor$ for "and" and "or", we typically use the multiplication and addition notation $\cdot\;+$ instead.

So, the previous example would be written as follow.

\[ x_1 x_2 + x_1 x_2 \]

In boolean satifiability, we also have the ability to set constraints such as $\mathrm{AtLeast}(n)$ and $\mathrm{AtMost}(n)$ to specify the least and the most amount of boolean variables that can be true respectively.

For example, by adding the constraint $\mathrm{AtLeast}(3)$ to the previous example, the solution $x_1 = \False$, $x_2 = \True$, $x_3 = \True$ would not be valid.

\newpage
\section*{Metric Dimension}

Let $G$ be a connected graph.

Let $V = \{v_1, v_2, \ldots, v_n\}$ be the set of all vertices in $G$.

Let $d(x, y)$ be the distance between vertices $x$ and $y$.

Let $W = \{w_1, w_2, \ldots, w_k\}$ be an ordered set of vertices in $G$.

For $v \in V$, the ordered tuple
	$$r(v|W) = (d(w_1, v), d(w_2, v), \ldots, d(w_k, v)$$
is the representation of $v$ with respect to $W$.

The set $W$ is a resolving set for $G$ if
	$$\forall u, v \in V, u \neq v \implies r(u|W) \neq r(v|W)$$
or in other words, all vertices in $G$ have distinct representations with respect to $W$.

The metric dimension of $G$ is the minimum cardinality of a resolving set for $G$.

\newpage
\section*{Metric Dimension Example}
\begin{minipage}[t]{0.65\textwidth}
\setlength{\parskip}{\baselineskip}

Let $W = \{x_2, x_4\}$.

\[
\begin{array}{lclcc}
r(x_1|W) &=& (\;d(x_2, x_1),\; d(x_4, x_1)\;) &=& \hl{$(\;2,\; 1\;)$} \\
r(x_2|W) &=& (\;d(x_2, x_2),\; d(x_4, x_2)\;) &=&      (\;0,\; 1\;)   \\
r(x_3|W) &=& (\;d(x_2, x_3),\; d(x_4, x_3)\;) &=& \hl{$(\;2,\; 1\;)$} \\
r(x_4|W) &=& (\;d(x_2, x_4),\; d(x_4, x_4)\;) &=&      (\;1,\; 0\;)   \\
r(x_5|W) &=& (\;d(x_2, x_5),\; d(x_4, x_5)\;) &=&      (\;1,\; 1\;)
\end{array}
\]

Therefore $\{x_2, x_4\}$ is not a resolving set.

\end{minipage}
\hfill
\begin{minipage}[t]{0.3\textwidth}
	\begin{figure}[H]
		


\def \globalscale {1.000000}
\begin{tikzpicture}[y=1cm, x=1cm, yscale=\globalscale,xscale=\globalscale, every node/.append style={scale=\globalscale}, inner sep=0pt, outer sep=0pt]
  \begin{scope}[shift={(0.141, -5.2283)}]
    \path[fill=white] (-0.141, 5.2283) -- (-0.141, 10.5975) -- (5.5035, 10.5975) -- (5.5035, 5.2283) -- (-0.141, 5.2283) -- (-0.141, 5.2283)-- cycle;;



    \path[draw=black] (0.9515, 6.5861) ellipse (0.9515cm and 0.6343cm);



    \node[anchor=south] (text9982) at (0.9515, 6.3951){x1};



    \path[draw=black] (4.411, 7.6443) ellipse (0.9515cm and 0.6343cm);



    \node[anchor=south] (text3045) at (4.411, 7.4533){x3};



    \path[draw=black] (1.8336, 6.856).. controls (2.3541, 7.0153) and (3.0106, 7.2162) .. (3.5307, 7.3751);



    \path[draw=black] (3.2742, 6.0035) ellipse (0.9515cm and 0.6343cm);



    \node[anchor=south] (text6271) at (3.2742, 5.8122){x4};



    \path[draw=black] (1.8561, 6.3591).. controls (2.0242, 6.3172) and (2.2001, 6.2731) .. (2.3682, 6.2308);



    \path[draw=black] (1.0674, 8.5137) ellipse (0.9515cm and 0.6343cm);



    \node[anchor=south] (text4006) at (1.0674, 8.3227){x5};



    \path[draw=black] (0.9899, 7.2281).. controls (1.0026, 7.4339) and (1.016, 7.662) .. (1.0287, 7.8681);



    \path[draw=black] (4.0072, 7.0615).. controls (3.8997, 6.9064) and (3.7841, 6.7394) .. (3.677, 6.5847);



    \path[draw=black] (3.5152, 7.8773).. controls (3.0353, 8.0024) and (2.446, 8.1553) .. (1.9657, 8.2804);



    \path[draw=black] (2.7843, 6.5607).. controls (2.4185, 6.9765) and (1.921, 7.5428) .. (1.5559, 7.9583);



    \path[draw=black] (3.3344, 9.8222) ellipse (0.9515cm and 0.6343cm);



    \node[anchor=south] (text5802) at (3.3344, 9.6308){x2};



    \path[draw=black] (3.3242, 9.1745).. controls (3.3133, 8.4721) and (3.2957, 7.3578) .. (3.2844, 6.6544);



    \path[draw=black] (2.6046, 9.4011).. controls (2.3459, 9.2516) and (2.0545, 9.0835) .. (1.7958, 8.9341);



  \end{scope}

\end{tikzpicture}

	\end{figure}
\end{minipage}

\newpage
\section*{Metric Dimension Example}
\begin{minipage}[t]{0.65\textwidth}
\setlength{\parskip}{\baselineskip}

Let $W = \{x_3, x_4\}$.

\[
\left[ \begin{array}{l}
	r(x_1|W) \\
	r(x_2|W) \\
	r(x_3|W) \\
	r(x_4|W) \\
	r(x_5|W)
\end{array} \right]
=
\left[ \begin{array}{ll}
	d(x_3,\; x_1) & d(x_4,\; x_1) \\
	d(x_3,\; x_2) & d(x_4,\; x_2) \\
	d(x_3,\; x_3) & d(x_4,\; x_3) \\
	d(x_3,\; x_4) & d(x_4,\; x_4) \\
	d(x_3,\; x_5) & d(x_4,\; x_5)
\end{array} \right]
=
\left[ \begin{array}{ll}
	\hlr 1 & 1 \\
	     2 & 1 \\
	     0 & 1 \\
	     1 & 0 \\
	\hlr 1 & 1
\end{array} \right]
\]

Therefore $\{x_3, x_4\}$ is not a resolving set.

\end{minipage}
\hfill
\begin{minipage}[t]{0.3\textwidth}
	\begin{figure}[H]
		


\def \globalscale {1.000000}
\begin{tikzpicture}[y=1cm, x=1cm, yscale=\globalscale,xscale=\globalscale, every node/.append style={scale=\globalscale}, inner sep=0pt, outer sep=0pt]
  \begin{scope}[shift={(0.141, -5.2283)}]
    \path[fill=white] (-0.141, 5.2283) -- (-0.141, 10.5975) -- (5.5035, 10.5975) -- (5.5035, 5.2283) -- (-0.141, 5.2283) -- (-0.141, 5.2283)-- cycle;;



    \path[draw=black] (0.9515, 6.5861) ellipse (0.9515cm and 0.6343cm);



    \node[anchor=south] (text9982) at (0.9515, 6.3951){x1};



    \path[draw=black] (4.411, 7.6443) ellipse (0.9515cm and 0.6343cm);



    \node[anchor=south] (text3045) at (4.411, 7.4533){x3};



    \path[draw=black] (1.8336, 6.856).. controls (2.3541, 7.0153) and (3.0106, 7.2162) .. (3.5307, 7.3751);



    \path[draw=black] (3.2742, 6.0035) ellipse (0.9515cm and 0.6343cm);



    \node[anchor=south] (text6271) at (3.2742, 5.8122){x4};



    \path[draw=black] (1.8561, 6.3591).. controls (2.0242, 6.3172) and (2.2001, 6.2731) .. (2.3682, 6.2308);



    \path[draw=black] (1.0674, 8.5137) ellipse (0.9515cm and 0.6343cm);



    \node[anchor=south] (text4006) at (1.0674, 8.3227){x5};



    \path[draw=black] (0.9899, 7.2281).. controls (1.0026, 7.4339) and (1.016, 7.662) .. (1.0287, 7.8681);



    \path[draw=black] (4.0072, 7.0615).. controls (3.8997, 6.9064) and (3.7841, 6.7394) .. (3.677, 6.5847);



    \path[draw=black] (3.5152, 7.8773).. controls (3.0353, 8.0024) and (2.446, 8.1553) .. (1.9657, 8.2804);



    \path[draw=black] (2.7843, 6.5607).. controls (2.4185, 6.9765) and (1.921, 7.5428) .. (1.5559, 7.9583);



    \path[draw=black] (3.3344, 9.8222) ellipse (0.9515cm and 0.6343cm);



    \node[anchor=south] (text5802) at (3.3344, 9.6308){x2};



    \path[draw=black] (3.3242, 9.1745).. controls (3.3133, 8.4721) and (3.2957, 7.3578) .. (3.2844, 6.6544);



    \path[draw=black] (2.6046, 9.4011).. controls (2.3459, 9.2516) and (2.0545, 9.0835) .. (1.7958, 8.9341);



  \end{scope}

\end{tikzpicture}

	\end{figure}
\end{minipage}

\newpage
\section*{Metric Dimension Example}
\begin{minipage}[t]{0.65\textwidth}
\setlength{\parskip}{\baselineskip}

In general, given the distance matrix $D$ of the graph, choose 2 columns for $W$ with cardinality 2. Then, check whether all rows are unique to determine if it forms a valid resolving set.

\[
\begin{array}{lcl}
D
&=&
\left[ \begin{array}{lllll}
	d(x_1,\;x_1)&d(x_2,\;x_1)&d(x_3,\;x_1)&d(x_4,\;x_1)&d(x_5,\;x_1)\\
	d(x_1,\;x_2)&d(x_2,\;x_2)&d(x_3,\;x_2)&d(x_4,\;x_2)&d(x_5,\;x_2)\\
	d(x_1,\;x_3)&d(x_2,\;x_3)&d(x_3,\;x_3)&d(x_4,\;x_3)&d(x_5,\;x_3)\\
	d(x_1,\;x_4)&d(x_2,\;x_4)&d(x_3,\;x_4)&d(x_4,\;x_4)&d(x_5,\;x_4)\\
	d(x_1,\;x_5)&d(x_2,\;x_5)&d(x_3,\;x_5)&d(x_4,\;x_5)&d(x_5,\;x_5)
\end{array} \right]
\\ \\
&=&
\left[ \begin{array}{lllll}
	0 & 2 & 1 & 1 & 1 \\
	2 & 0 & 2 & 1 & 1 \\
	1 & 2 & 0 & 1 & 1 \\
	1 & 1 & 1 & 0 & 1 \\
	1 & 1 & 1 & 1 & 0
\end{array} \right]
\end{array}
\]

\end{minipage}
\hfill
\begin{minipage}[t]{0.3\textwidth}
	\begin{figure}[H]
		


\def \globalscale {1.000000}
\begin{tikzpicture}[y=1cm, x=1cm, yscale=\globalscale,xscale=\globalscale, every node/.append style={scale=\globalscale}, inner sep=0pt, outer sep=0pt]
  \begin{scope}[shift={(0.141, -5.2283)}]
    \path[fill=white] (-0.141, 5.2283) -- (-0.141, 10.5975) -- (5.5035, 10.5975) -- (5.5035, 5.2283) -- (-0.141, 5.2283) -- (-0.141, 5.2283)-- cycle;;



    \path[draw=black] (0.9515, 6.5861) ellipse (0.9515cm and 0.6343cm);



    \node[anchor=south] (text9982) at (0.9515, 6.3951){x1};



    \path[draw=black] (4.411, 7.6443) ellipse (0.9515cm and 0.6343cm);



    \node[anchor=south] (text3045) at (4.411, 7.4533){x3};



    \path[draw=black] (1.8336, 6.856).. controls (2.3541, 7.0153) and (3.0106, 7.2162) .. (3.5307, 7.3751);



    \path[draw=black] (3.2742, 6.0035) ellipse (0.9515cm and 0.6343cm);



    \node[anchor=south] (text6271) at (3.2742, 5.8122){x4};



    \path[draw=black] (1.8561, 6.3591).. controls (2.0242, 6.3172) and (2.2001, 6.2731) .. (2.3682, 6.2308);



    \path[draw=black] (1.0674, 8.5137) ellipse (0.9515cm and 0.6343cm);



    \node[anchor=south] (text4006) at (1.0674, 8.3227){x5};



    \path[draw=black] (0.9899, 7.2281).. controls (1.0026, 7.4339) and (1.016, 7.662) .. (1.0287, 7.8681);



    \path[draw=black] (4.0072, 7.0615).. controls (3.8997, 6.9064) and (3.7841, 6.7394) .. (3.677, 6.5847);



    \path[draw=black] (3.5152, 7.8773).. controls (3.0353, 8.0024) and (2.446, 8.1553) .. (1.9657, 8.2804);



    \path[draw=black] (2.7843, 6.5607).. controls (2.4185, 6.9765) and (1.921, 7.5428) .. (1.5559, 7.9583);



    \path[draw=black] (3.3344, 9.8222) ellipse (0.9515cm and 0.6343cm);



    \node[anchor=south] (text5802) at (3.3344, 9.6308){x2};



    \path[draw=black] (3.3242, 9.1745).. controls (3.3133, 8.4721) and (3.2957, 7.3578) .. (3.2844, 6.6544);



    \path[draw=black] (2.6046, 9.4011).. controls (2.3459, 9.2516) and (2.0545, 9.0835) .. (1.7958, 8.9341);



  \end{scope}

\end{tikzpicture}

	\end{figure}
\end{minipage}



\newpage
\section*{Metric Dimension Example}
\begin{minipage}[t]{0.65\textwidth}
\setlength{\parskip}{\baselineskip}

In this example, there are no valid resolving set with cardinality 2.

\[
\left[ \begin{array}{ll}
	0 & 2 \\
	2 & 0 \\
	1 & 2 \\
	1 & 1 \\
	1 & 1
\end{array} \right]
\quad,\quad
\left[ \begin{array}{ll}
	0 & 1 \\
	2 & 2 \\
	1 & 0 \\
	1 & 1 \\
	1 & 1
\end{array} \right]
\quad,\quad
\left[ \begin{array}{ll}
	0 & 1 \\
	2 & 1 \\
	1 & 1 \\
	1 & 0 \\
	1 & 1
\end{array} \right]
\quad,\quad
\left[ \begin{array}{ll}
	0 & 1 \\
	2 & 1 \\
	1 & 1 \\
	1 & 1 \\
	1 & 0
\end{array} \right]
\quad,\quad
\cdots
\]

\end{minipage}
\hfill
\begin{minipage}[t]{0.3\textwidth}
	\begin{figure}[H]
		


\def \globalscale {1.000000}
\begin{tikzpicture}[y=1cm, x=1cm, yscale=\globalscale,xscale=\globalscale, every node/.append style={scale=\globalscale}, inner sep=0pt, outer sep=0pt]
  \begin{scope}[shift={(0.141, -5.2283)}]
    \path[fill=white] (-0.141, 5.2283) -- (-0.141, 10.5975) -- (5.5035, 10.5975) -- (5.5035, 5.2283) -- (-0.141, 5.2283) -- (-0.141, 5.2283)-- cycle;;



    \path[draw=black] (0.9515, 6.5861) ellipse (0.9515cm and 0.6343cm);



    \node[anchor=south] (text9982) at (0.9515, 6.3951){x1};



    \path[draw=black] (4.411, 7.6443) ellipse (0.9515cm and 0.6343cm);



    \node[anchor=south] (text3045) at (4.411, 7.4533){x3};



    \path[draw=black] (1.8336, 6.856).. controls (2.3541, 7.0153) and (3.0106, 7.2162) .. (3.5307, 7.3751);



    \path[draw=black] (3.2742, 6.0035) ellipse (0.9515cm and 0.6343cm);



    \node[anchor=south] (text6271) at (3.2742, 5.8122){x4};



    \path[draw=black] (1.8561, 6.3591).. controls (2.0242, 6.3172) and (2.2001, 6.2731) .. (2.3682, 6.2308);



    \path[draw=black] (1.0674, 8.5137) ellipse (0.9515cm and 0.6343cm);



    \node[anchor=south] (text4006) at (1.0674, 8.3227){x5};



    \path[draw=black] (0.9899, 7.2281).. controls (1.0026, 7.4339) and (1.016, 7.662) .. (1.0287, 7.8681);



    \path[draw=black] (4.0072, 7.0615).. controls (3.8997, 6.9064) and (3.7841, 6.7394) .. (3.677, 6.5847);



    \path[draw=black] (3.5152, 7.8773).. controls (3.0353, 8.0024) and (2.446, 8.1553) .. (1.9657, 8.2804);



    \path[draw=black] (2.7843, 6.5607).. controls (2.4185, 6.9765) and (1.921, 7.5428) .. (1.5559, 7.9583);



    \path[draw=black] (3.3344, 9.8222) ellipse (0.9515cm and 0.6343cm);



    \node[anchor=south] (text5802) at (3.3344, 9.6308){x2};



    \path[draw=black] (3.3242, 9.1745).. controls (3.3133, 8.4721) and (3.2957, 7.3578) .. (3.2844, 6.6544);



    \path[draw=black] (2.6046, 9.4011).. controls (2.3459, 9.2516) and (2.0545, 9.0835) .. (1.7958, 8.9341);



  \end{scope}

\end{tikzpicture}

	\end{figure}
\end{minipage}

\newpage
\section*{Metric Dimension Example}
\begin{minipage}[t]{0.65\textwidth}
\setlength{\parskip}{\baselineskip}

The metric dimension of this graph is 3. One valid resolving set for the graph is $\{x_1, x_4, x_5\}$, another is $\{x_1, x_3, x_4\}$, etc.

\[
\left[ \begin{array}{>{\columncolor{yellow}}lll>{\columncolor{yellow}}l>{\columncolor{yellow}}l}
	0 & 2 & 1 & 1 & 1 \\
	2 & 0 & 2 & 1 & 1 \\
	1 & 2 & 0 & 1 & 1 \\
	1 & 1 & 1 & 0 & 1 \\
	1 & 1 & 1 & 1 & 0
\end{array} \right]
\quad,\quad
\left[ \begin{array}{>{\columncolor{yellow}}ll>{\columncolor{yellow}}l>{\columncolor{yellow}}ll}
	0 & 2 & 1 & 1 & 1 \\
	2 & 0 & 2 & 1 & 1 \\
	1 & 2 & 0 & 1 & 1 \\
	1 & 1 & 1 & 0 & 1 \\
	1 & 1 & 1 & 1 & 0
\end{array} \right]
\quad,\quad
\cdots
\]

There are 8 valid resolving sets for this graph.

\end{minipage}
\hfill
\begin{minipage}[t]{0.3\textwidth}
	\begin{figure}[H]
		


\def \globalscale {1.000000}
\begin{tikzpicture}[y=1cm, x=1cm, yscale=\globalscale,xscale=\globalscale, every node/.append style={scale=\globalscale}, inner sep=0pt, outer sep=0pt]
  \begin{scope}[shift={(0.141, -5.2283)}]
    \path[fill=white] (-0.141, 5.2283) -- (-0.141, 10.5975) -- (5.5035, 10.5975) -- (5.5035, 5.2283) -- (-0.141, 5.2283) -- (-0.141, 5.2283)-- cycle;;



    \path[draw=black] (0.9515, 6.5861) ellipse (0.9515cm and 0.6343cm);



    \node[anchor=south] (text9982) at (0.9515, 6.3951){x1};



    \path[draw=black] (4.411, 7.6443) ellipse (0.9515cm and 0.6343cm);



    \node[anchor=south] (text3045) at (4.411, 7.4533){x3};



    \path[draw=black] (1.8336, 6.856).. controls (2.3541, 7.0153) and (3.0106, 7.2162) .. (3.5307, 7.3751);



    \path[draw=black] (3.2742, 6.0035) ellipse (0.9515cm and 0.6343cm);



    \node[anchor=south] (text6271) at (3.2742, 5.8122){x4};



    \path[draw=black] (1.8561, 6.3591).. controls (2.0242, 6.3172) and (2.2001, 6.2731) .. (2.3682, 6.2308);



    \path[draw=black] (1.0674, 8.5137) ellipse (0.9515cm and 0.6343cm);



    \node[anchor=south] (text4006) at (1.0674, 8.3227){x5};



    \path[draw=black] (0.9899, 7.2281).. controls (1.0026, 7.4339) and (1.016, 7.662) .. (1.0287, 7.8681);



    \path[draw=black] (4.0072, 7.0615).. controls (3.8997, 6.9064) and (3.7841, 6.7394) .. (3.677, 6.5847);



    \path[draw=black] (3.5152, 7.8773).. controls (3.0353, 8.0024) and (2.446, 8.1553) .. (1.9657, 8.2804);



    \path[draw=black] (2.7843, 6.5607).. controls (2.4185, 6.9765) and (1.921, 7.5428) .. (1.5559, 7.9583);



    \path[draw=black] (3.3344, 9.8222) ellipse (0.9515cm and 0.6343cm);



    \node[anchor=south] (text5802) at (3.3344, 9.6308){x2};



    \path[draw=black] (3.3242, 9.1745).. controls (3.3133, 8.4721) and (3.2957, 7.3578) .. (3.2844, 6.6544);



    \path[draw=black] (2.6046, 9.4011).. controls (2.3459, 9.2516) and (2.0545, 9.0835) .. (1.7958, 8.9341);



  \end{scope}

\end{tikzpicture}

	\end{figure}
\end{minipage}

\newpage
\section*{Metric Dimension in Boolean Satisfiability}

Here, our objective is to translate the problem of choosing columns from the distance matrix and testing for any duplicate rows into something that is computable by boolean satisfiability.

First, we compute all the pair-wise combination of the rows for equality.

\[
\left[ \begin{array}{lllll}
	\hlr 0 & 2 & 1 & 1 & 1 \\
	     2 & 0 & 2 & 1 & 1 \\
	     1 & 2 & 0 & 1 & 1 \\
	     1 & 1 & 1 & 0 & 1 \\
	     1 & 1 & 1 & 1 & 0
\end{array} \right]
\mapsto
\left[ \begin{array}{ccccc}
	\\
	\False & \False & \False &  \True &  \True \\
	\False &  \True & \False &  \True &  \True \\
	\False & \False &  \True & \False &  \True \\
	\False & \False &  \True &  \True & \False
\end{array} \right]
\]

\[
\left[ \begin{array}{lllll}
	     0 & 2 & 1 & 1 & 1 \\
	\hlr 2 & 0 & 2 & 1 & 1 \\
	     1 & 2 & 0 & 1 & 1 \\
	     1 & 1 & 1 & 0 & 1 \\
	     1 & 1 & 1 & 1 & 0
\end{array} \right]
\mapsto
\left[ \begin{array}{ccccc}
	\\
	\\
	\False & \False & \False &  \True &  \True \\
	\False & \False & \False & \False &  \True \\
	\False & \False & \False &  \True & \False
\end{array} \right]
\]

\[
\left[ \begin{array}{lllll}
	     0 & 2 & 1 & 1 & 1 \\
	     2 & 0 & 2 & 1 & 1 \\
	\hlr 1 & 2 & 0 & 1 & 1 \\
	     1 & 1 & 1 & 0 & 1 \\
	     1 & 1 & 1 & 1 & 0
\end{array} \right]
\mapsto
\left[ \begin{array}{ccccc}
	\\
	\\
	\\
	 \True & \False & \False & \False &  \True \\
	 \True & \False & \False &  \True & \False \\
\end{array} \right]
\]

\newpage
\section*{Metric Dimension in Boolean Satisfiability}

\[
\left[ \begin{array}{lllll}
	     0 & 2 & 1 & 1 & 1 \\
	     2 & 0 & 2 & 1 & 1 \\
	     1 & 2 & 0 & 1 & 1 \\
	\hlr 1 & 1 & 1 & 0 & 1 \\
	     1 & 1 & 1 & 1 & 0
\end{array} \right]
\mapsto
\left[ \begin{array}{ccccc}
	\\
	\\
	\\
	\\
	 \True &  \True &  \True & \False & \False
\end{array} \right]
\]

Then combine the results into one matrix, removing duplicates.

\begin{minipage}[m]{0.55\textwidth}
\setlength{\parskip}{\baselineskip}

\[
\left[ \begin{array}{ccccc}
	\False & \False & \False & \False &  \True \\
	\False & \False & \False &  \True & \False \\
	\False & \False & \False &  \True &  \True \\
	\False & \False &  \True & \False &  \True \\
	\False & \False &  \True &  \True & \False \\
	\False &  \True & \False &  \True &  \True \\
	 \True & \False & \False & \False &  \True \\
	 \True & \False & \False &  \True & \False \\
	 \True &  \True &  \True & \False & \False
\end{array} \right]
\]

\end{minipage}
\hfill
\begin{minipage}[m]{0.4\textwidth}
	\begin{figure}[H]
		


\def \globalscale {1.000000}
\begin{tikzpicture}[y=1cm, x=1cm, yscale=\globalscale,xscale=\globalscale, every node/.append style={scale=\globalscale}, inner sep=0pt, outer sep=0pt]
  \begin{scope}[shift={(0.141, -5.2283)}]
    \path[fill=white] (-0.141, 5.2283) -- (-0.141, 10.5975) -- (5.5035, 10.5975) -- (5.5035, 5.2283) -- (-0.141, 5.2283) -- (-0.141, 5.2283)-- cycle;;



    \path[draw=black] (0.9515, 6.5861) ellipse (0.9515cm and 0.6343cm);



    \node[anchor=south] (text9982) at (0.9515, 6.3951){x1};



    \path[draw=black] (4.411, 7.6443) ellipse (0.9515cm and 0.6343cm);



    \node[anchor=south] (text3045) at (4.411, 7.4533){x3};



    \path[draw=black] (1.8336, 6.856).. controls (2.3541, 7.0153) and (3.0106, 7.2162) .. (3.5307, 7.3751);



    \path[draw=black] (3.2742, 6.0035) ellipse (0.9515cm and 0.6343cm);



    \node[anchor=south] (text6271) at (3.2742, 5.8122){x4};



    \path[draw=black] (1.8561, 6.3591).. controls (2.0242, 6.3172) and (2.2001, 6.2731) .. (2.3682, 6.2308);



    \path[draw=black] (1.0674, 8.5137) ellipse (0.9515cm and 0.6343cm);



    \node[anchor=south] (text4006) at (1.0674, 8.3227){x5};



    \path[draw=black] (0.9899, 7.2281).. controls (1.0026, 7.4339) and (1.016, 7.662) .. (1.0287, 7.8681);



    \path[draw=black] (4.0072, 7.0615).. controls (3.8997, 6.9064) and (3.7841, 6.7394) .. (3.677, 6.5847);



    \path[draw=black] (3.5152, 7.8773).. controls (3.0353, 8.0024) and (2.446, 8.1553) .. (1.9657, 8.2804);



    \path[draw=black] (2.7843, 6.5607).. controls (2.4185, 6.9765) and (1.921, 7.5428) .. (1.5559, 7.9583);



    \path[draw=black] (3.3344, 9.8222) ellipse (0.9515cm and 0.6343cm);



    \node[anchor=south] (text5802) at (3.3344, 9.6308){x2};



    \path[draw=black] (3.3242, 9.1745).. controls (3.3133, 8.4721) and (3.2957, 7.3578) .. (3.2844, 6.6544);



    \path[draw=black] (2.6046, 9.4011).. controls (2.3459, 9.2516) and (2.0545, 9.0835) .. (1.7958, 8.9341);



  \end{scope}

\end{tikzpicture}

	\end{figure}
\end{minipage}

\newpage
\section*{Metric Dimension in Boolean Satisfiability}

Here, this distance similarity matrix can be interpreted as showing where the invalid resolving are.

\[
\begin{array}{ccccccl}
	\False & \False & \False & \False &  \True & \mapsto & \{x_5\} \\
	\False & \False & \False &  \True & \False & \mapsto & \{x_4\} \\
	\False & \False & \False &  \True &  \True & \mapsto & \{x_4, x_5\} \\
	\False & \False &  \True & \False &  \True & \mapsto & \{x_3, x_5\} \\
	\False & \False &  \True &  \True & \False & \mapsto & \{x_3, x_4\} \\
	\False &  \True & \False &  \True &  \True & \mapsto & \{x_2, x_4, x_5\} \\
	 \True & \False & \False & \False &  \True & \mapsto & \{x_1, x_5\} \\
	 \True & \False & \False &  \True & \False & \mapsto & \{x_1, x_4\} \\
	 \True &  \True &  \True & \False & \False & \mapsto & \{x_1, x_2, x_3\}
\end{array}
\]

\newpage
\section*{Metric Dimension in Boolean Satisfiability}

However, this interpretation is incomplete. Consider the following, for example.

\[
\begin{array}{ccccccl}
	\False &  \True & \False &  \True &  \True & \mapsto & \{x_2, x_4, x_5\} \\
\end{array}
\]

It is also true to interpret the following sets as invalid resolving set.

\[ \{x_2\},\; \{x_4\},\; \{x_5\},\; \{x_2, x_4\},\; \{x_2, x_5\},\; \{x_4, x_5\} \]

As such, after taking the row equality, it should be expanded like follow.

\[
\begin{array}{ccccccl}
	\False &  \True & \False &  \True &  \True & \mapsto & \{x_2, x_4, x_5\} \\
	\False &  \True & \False &  \True & \False & \mapsto & \{x_2, x_4\} \\
	\False &  \True & \False & \False &  \True & \mapsto & \{x_2, x_5\} \\
	\False & \False & \False &  \True &  \True & \mapsto & \{x_4, x_5\} \\
	\False &  \True & \False & \False & \False & \mapsto & \{x_2\} \\
	\False & \False & \False &  \True & \False & \mapsto & \{x_4\} \\
	\False & \False & \False & \False &  \True & \mapsto & \{x_5\}
\end{array}
\]

\newpage
\section*{Metric Dimension in Boolean Satisfiability}

So, the second step is to do this permuting operation for every single result in the distance similarity matrix. Then, combine all of the result back into one matrix, removing duplicates.

\[
\begin{array}{ccccccl}
	\False & \False & \False & \False &  \True & \mapsto & \{x_5\} \\
	\False & \False & \False &  \True & \False & \mapsto & \{x_4\} \\
	\False & \False & \False &  \True &  \True & \mapsto & \{x_4, x_5\} \\
	\False & \False &  \True & \False & \False & \mapsto & \{x_3\} \\
	\False & \False &  \True & \False &  \True & \mapsto & \{x_3, x_5\} \\
	\False & \False &  \True &  \True & \False & \mapsto & \{x_3, x_4\} \\
	\False &  \True & \False & \False & \False & \mapsto & \{x_2\} \\
	\False &  \True & \False & \False &  \True & \mapsto & \{x_2, x_5\} \\
	\False &  \True & \False &  \True & \False & \mapsto & \{x_2, x_4\} \\
	\False &  \True & \False &  \True &  \True & \mapsto & \{x_2, x_4, x_5\} \\
	\False &  \True &  \True & \False & \False & \mapsto & \{x_2, x_3\} \\
	 \True & \False & \False & \False & \False & \mapsto & \{x_1\} \\
	 \True & \False & \False & \False &  \True & \mapsto & \{x_1, x_5\} \\
	 \True & \False & \False &  \True & \False & \mapsto & \{x_1, x_4\} \\
	 \True & \False &  \True & \False & \False & \mapsto & \{x_1, x_3\} \\
	 \True &  \True & \False & \False & \False & \mapsto & \{x_1, x_2\} \\
	 \True &  \True &  \True & \False & \False & \mapsto & \{x_1, x_2, x_3\}
\end{array}
\]

\newpage
\section*{Metric Dimension in Boolean Satisfiability}

Third, from our permuted distance similarity matrix, we can immediately translate it into a boolean satisfiability expression that states the various invalid resolving set.

\begin{align*}
	\neg ( x_5 + x_4 &+ x_4 x_5 + x_3 + x_3 x_5 + x_3 x_4 + x_2 + x_2 x_5 + x_2 x_4 \\ &+ x_2 x_4 x_5 + x_2 x_3 + x_1 + x_1 x_5 + x_1 x_4 + x_1 x_3 + x_1 x_2 + x_1 x_2 x_3 )
\end{align*}

Evaluating this will result in a model that is a valid resolving set.

Fourth, to find the minimum cardinality, simply use $\mathrm{AtLeast}$ and $\mathrm{AtMost}$ to set the cardinality, then iterate over from 1 and onwards, taking the first cardinality that is satisfiable as the metric dimension.

\end{document}
